% Part: first-order-logic
% Chapter: completeness
% Section: complete-sets

% Définition des ensembles cohérents complets. Les propriétés des
% ensembles complets nécessaires à la complétude sont établies ici à
% partir des résultats de dérivabilité du système de preuve choisi.

\documentclass[../../../include/open-logic-section]{subfiles}

\begin{document}

\iftag{FOL}
      {\olfileid[fr]{fol}{com}{ccs}}
      {\olfileid[fr]{pl}{com}{ccs}}

\Needspace{7\baselineskip}
\iftag{FOL}
{\olsection{Ensembles cohérents et complets d'\usetoken{p}{sentence}}}
{\olsection{Ensembles cohérents et complets de formules}}

\begin{defn}[Ensemble complet]
\ollabel{def:complete-set} Un ensemble~$\Gamma$
\iftag{FOL}{d'!!{sentence}s}{de !!{formula}s} est
\emph{!!{complete}} si et seulement si, pour
\iftag{FOL}{tout !!{sentence}}{toute !!{formula}}~$!A$,
on a $!A\in\Gamma$ ou $\lnot !A\in\Gamma$.
\end{defn}

\begin{explain}
Les ensembles !!{complete}s \iftag{FOL}{d'!!{sentence}s}{de !!{formula}s} ne laissent aucune
question en suspens : pour chaque \iftag{FOL}{!!{sentence}}{!!{formula}}~$!A$, $\Gamma$
« affirme » $!A$ ou sa négation. Si l'on exige aussi la cohérence, il
ne peut pas affirmer les deux. L'intérêt des ensembles !!{complete}s
dépasse la démonstration du théorème de complétude. Par exemple, une
théorie cohérente, complète et effectivement axiomatisable est
décidable. On énumère les théorèmes à partir de l'axiomatisation
effective jusqu'à rencontrer $!A$ ou $\lnot !A$. La complétude
garantit que l'énumération atteindra l'un des deux, et la cohérence
garantit l'unicité de la réponse.
\end{explain}

\begin{explain}
Les ensembles cohérents et !!{complete}s jouent un rôle important
dans la démonstration de complétude : nous pouvons garantir que tout
ensemble cohérent \iftag{FOL}{d'!!{sentence}s}{de !!{formula}s}~$\Gamma$ est inclus dans un
ensemble cohérent et !!{complete}~$\Gamma^*$. Celui-ci contient,
pour chaque \iftag{FOL}{!!{sentence}}{!!{formula}}~$!A$, soit $!A$, soit sa négation
$\lnot !A$, mais non les deux. C'est notamment vrai pour les
\iftag{FOL}{!!{sentence}s atomiques}{!!{propositional variable}s}.
À partir d'un ensemble cohérent et !!{complete}\iftag{FOL}{,
dans un langage convenablement enrichi en !!{constant}s}{}, nous pouvons donc construire
\iftag{FOL}{!!a{structure}}{!!a{valuation}} dont
\iftag{FOL}{l'interprétation des !!{predicate}s}{la valeur de vérité
attribuée aux variables propositionnelles} est définie d'après les
\iftag{FOL}{!!{sentence}s atomiques}{!!{propositional variable}s}
appartenant à~$\Gamma^*$. Nous montrerons ensuite que cette
\iftag{FOL}{!!{structure}}{!!{valuation}} satisfait \iftag{FOL}{tous les !!{sentence}s}{toutes les !!{formula}s}
de~$\Gamma^*$, et donc aussi chaque élément de~$\Gamma$.
Pour cela, il nous faut notamment les équivalences
$\lnot !A\in\Gamma^*$ si et seulement si
$!A\notin\Gamma^*$, et $(!A\lor !B)\in\Gamma^*$ si et seulement
si $!A\in\Gamma^*$ ou $!B\in\Gamma^*$.
\end{explain}

Dans la suite, nous utiliserons souvent sans les rappeler la
réflexivité, la monotonie et la transitivité de~$\Proves$ (voir les
\iftag{FOL}{%
  \tagrefs{prfSC/{fol:seq:ptn:sec},prfND/{fol:ntd:ptn:sec},prfAX/{fol:axd:ptn:sec},prfTab/{fol:tab:ptn:sec}}}{%
  \tagrefs{prfSC/{pl:seq:ptn:sec},prfND/{pl:ntd:ptn:sec},prfAX/{pl:axd:ptn:sec},prfTab/{pl:tab:ptn:sec}}}).

\begin{prop}
\ollabel{prop:ccs}
Supposons $\Gamma$ cohérent et !!{complete}. Pour
\iftag{FOL}{des !!{sentence}s}{des !!{formula}s} $!A$ et~$!B$,
on a les propriétés suivantes.
\begin{enumerate}
\item \ollabel{prop:ccs-prov-in} Si $\Gamma\Proves !A$, alors
  $!A\in\Gamma$.

\tagitem{prvAnd}{\ollabel{prop:ccs-and} $!A\land !B\in\Gamma$
  si et seulement si $!A\in\Gamma$ et $!B\in\Gamma$.}{}

\tagitem{prvOr}{\ollabel{prop:ccs-or} $!A\lor !B\in\Gamma$ si et
  seulement si $!A\in\Gamma$ ou $!B\in\Gamma$.}{}

\tagitem{prvIf}{\ollabel{prop:ccs-if} $!A\lif !B\in\Gamma$ si et
  seulement si $!A\notin\Gamma$ ou $!B\in\Gamma$.}{}
\end{enumerate}
\end{prop}

\begin{proof}
Supposons, pour toute la démonstration, que $\Gamma$ soit cohérent
et !!{complete}.
\begin{enumerate}
\item Si $\Gamma\Proves !A$, alors $!A\in\Gamma$.

Supposons $\Gamma\Proves !A$ et, par l'absurde,
$!A\notin\Gamma$. Comme $\Gamma$ est !!{complete},
$\lnot !A\in\Gamma$. D'après les
\iftag{FOL}{%
  \tagrefs{prfSC/{fol:seq:prv:prop:explicit-inc},
    prfND/{fol:ntd:prv:prop:explicit-inc},
    prfAX/{fol:axd:prv:prop:explicit-inc},
    prfTab/{fol:tab:prv:prop:explicit-inc}}}{%
  \tagrefs{prfSC/{pl:seq:prv:prop:explicit-inc},
    prfND/{pl:ntd:prv:prop:explicit-inc},
    prfAX/{pl:axd:prv:prop:explicit-inc},
    prfTab/{pl:tab:prv:prop:explicit-inc}}},
$\Gamma$ est incohérent, contrairement à l'hypothèse. Ainsi,
$!A\notin\Gamma$ est impossible et $!A\in\Gamma$.

\tagitem{defAnd}{}{%
\iftag{probAnd}{Exercice.}{%
$!A\land !B\in\Gamma$ si et seulement si $!A\in\Gamma$ et
$!B\in\Gamma$.

Pour le sens direct, supposons $!A\land !B\in\Gamma$. Par
réflexivité, $\Gamma\Proves !A\land !B$. D'après les
\iftag{FOL}{%
  \tagrefs{prfSC/{fol:seq:ppr:prop:provability-land},
    prfND/{fol:ntd:ppr:prop:provability-land},
    prfAX/{fol:axd:ppr:prop:provability-land},
    prfTab/{fol:tab:ppr:prop:provability-land}}}{%
  \tagrefs{prfSC/{pl:seq:ppr:prop:provability-land},
    prfND/{pl:ntd:ppr:prop:provability-land},
    prfAX/{pl:axd:ppr:prop:provability-land},
    prfTab/{pl:tab:ppr:prop:provability-land}}}, point~(1),
$\Gamma\Proves !A$ et $\Gamma\Proves !B$. Par
\olref{prop:ccs}\olref{prop:ccs-prov-in}, $!A\in\Gamma$ et $!B\in\Gamma$.

Réciproquement, supposons $!A\in\Gamma$ et $!B\in\Gamma$.
Par réflexivité, $\Gamma\Proves !A$ et $\Gamma\Proves !B$. D'après les
\iftag{FOL}{%
  \tagrefs{prfSC/{fol:seq:ppr:prop:provability-land},
    prfND/{fol:ntd:ppr:prop:provability-land},
    prfAX/{fol:axd:ppr:prop:provability-land},
    prfTab/{fol:tab:ppr:prop:provability-land}}}{%
  \tagrefs{prfSC/{pl:seq:ppr:prop:provability-land},
    prfND/{pl:ntd:ppr:prop:provability-land},
    prfAX/{pl:axd:ppr:prop:provability-land},
    prfTab/{pl:tab:ppr:prop:provability-land}}}, point~(2),
$\Gamma\Proves !A\land !B$. La
\olref{prop:ccs}\olref{prop:ccs-prov-in} donne $!A\land !B\in\Gamma$.}}

\tagitem{defOr}{}{%
\iftag{probOr}{Exercice.}{%
Montrons d'abord que $!A\lor !B\in\Gamma$ entraîne
$!A\in\Gamma$ ou $!B\in\Gamma$. Supposons au contraire que
$!A\lor !B\in\Gamma$, mais $!A\notin\Gamma$ et
$!B\notin\Gamma$. La complétude de $\Gamma$ donne
$\lnot !A\in\Gamma$ et $\lnot !B\in\Gamma$. D'après les
\iftag{FOL}{%
  \tagrefs{prfSC/{fol:seq:ppr:prop:provability-lor},
    prfND/{fol:ntd:ppr:prop:provability-lor},
    prfAX/{fol:axd:ppr:prop:provability-lor},
    prfTab/{fol:tab:ppr:prop:provability-lor}}}{%
  \tagrefs{prfSC/{pl:seq:ppr:prop:provability-lor},
    prfND/{pl:ntd:ppr:prop:provability-lor},
    prfAX/{pl:axd:ppr:prop:provability-lor},
    prfTab/{pl:tab:ppr:prop:provability-lor}}}, point~(1),
$\Gamma$ serait incohérent. C'est contradictoire ; donc
$!A\in\Gamma$ ou $!B\in\Gamma$.

Pour la réciproque, supposons $!A\in\Gamma$ ou
$!B\in\Gamma$. Par réflexivité et d'après les
\iftag{FOL}{%
  \tagrefs{prfSC/{fol:seq:ppr:prop:provability-lor},
    prfND/{fol:ntd:ppr:prop:provability-lor},
    prfAX/{fol:axd:ppr:prop:provability-lor},
    prfTab/{fol:tab:ppr:prop:provability-lor}}}{%
  \tagrefs{prfSC/{pl:seq:ppr:prop:provability-lor},
    prfND/{pl:ntd:ppr:prop:provability-lor},
    prfAX/{pl:axd:ppr:prop:provability-lor},
    prfTab/{pl:tab:ppr:prop:provability-lor}}}, point~(2),
$\Gamma\Proves !A\lor !B$. La
\olref{prop:ccs}\olref{prop:ccs-prov-in} donne $!A\lor !B\in\Gamma$.}}

\tagitem{defIf}{}{%
\iftag{probIf}{Exercice.}{%
Pour le sens direct, supposons $!A\lif !B\in\Gamma$, puis,
par l'absurde, $!A\in\Gamma$ et $!B\notin\Gamma$. Les deux
premières appartenances donnent
$\Gamma\Proves !A\lif !B$ et $\Gamma\Proves !A$ par
réflexivité. D'après les
\iftag{FOL}{%
  \tagrefs{prfSC/{fol:seq:ppr:prop:provability-lif},
    prfND/{fol:ntd:ppr:prop:provability-lif},
    prfAX/{fol:axd:ppr:prop:provability-lif},
    prfTab/{fol:tab:ppr:prop:provability-lif}}}{%
  \tagrefs{prfSC/{pl:seq:ppr:prop:provability-lif},
    prfND/{pl:ntd:ppr:prop:provability-lif},
    prfAX/{pl:axd:ppr:prop:provability-lif},
    prfTab/{pl:tab:ppr:prop:provability-lif}}}, point~(1),
$\Gamma\Proves !B$. La
\olref{prop:ccs}\olref{prop:ccs-prov-in} donne $!B\in\Gamma$, contrairement à
l'hypothèse.

Pour la réciproque, considérons d'abord le cas
$!A\notin\Gamma$. Puisque $\Gamma$ est !!{complete},
$\lnot !A\in\Gamma$. D'après les
\iftag{FOL}{%
  \tagrefs{prfSC/{fol:seq:ppr:prop:provability-lif},
    prfND/{fol:ntd:ppr:prop:provability-lif},
    prfAX/{fol:axd:ppr:prop:provability-lif},
    prfTab/{fol:tab:ppr:prop:provability-lif}}}{%
  \tagrefs{prfSC/{pl:seq:ppr:prop:provability-lif},
    prfND/{pl:ntd:ppr:prop:provability-lif},
    prfAX/{pl:axd:ppr:prop:provability-lif},
    prfTab/{pl:tab:ppr:prop:provability-lif}}}, point~(2),
$\Gamma\Proves !A\lif !B$. La
\olref{prop:ccs}\olref{prop:ccs-prov-in} donne $!A\lif !B\in\Gamma$.

Considérons maintenant le cas $!B\in\Gamma$. D'après les
\iftag{FOL}{%
  \tagrefs{prfSC/{fol:seq:ppr:prop:provability-lif},
    prfND/{fol:ntd:ppr:prop:provability-lif},
    prfTab/{fol:tab:ppr:prop:provability-lif},
    prfAX/{fol:axd:ppr:prop:provability-lif}}}{%
  \tagrefs{prfSC/{pl:seq:ppr:prop:provability-lif},
    prfND/{pl:ntd:ppr:prop:provability-lif},
    prfTab/{pl:tab:ppr:prop:provability-lif},
    prfAX/{pl:axd:ppr:prop:provability-lif}}},
point~(2), $\Gamma\Proves !A\lif !B$. La
\olref{prop:ccs}\olref{prop:ccs-prov-in} donne $!A\lif !B\in\Gamma$.}}
\end{enumerate}
\end{proof}

\tagprob{FOL}
\begin{prob}
Achever la démonstration de la \olref[fol][com][ccs]{prop:ccs}.
\end{prob}
\tagendprob

\tagprob{notFOL}
\begin{prob}
Achever la démonstration de la \olref[pl][com][ccs]{prop:ccs}.
\end{prob}
\tagendprob

\end{document}
